\renewcommand{\tema}{Preguntas Post-Clase - Ondas}

\twocolumn
\section*{PREGUNTAS POST-CLASE}

% -------------------------------------------------------
%  PREGUNTAS 1-5: Aplicación directa de conceptos
% -------------------------------------------------------

\begin{preguntabox}
\section{Pregunta 1}
Un submarino utiliza sonar para detectar objetos bajo el agua. Las ondas de sonar son:

\begin{enumerate}
    \item[A)] Transversales, porque viajan en un líquido.
    \item[B)] Longitudinales, porque son ondas de presión que se comprimen y rarifican en la dirección de propagación.
    \item[C)] Transversales, porque la perturbación es perpendicular al movimiento del submarino.
    \item[D)] Electromagnéticas, porque no requieren medio material.
\end{enumerate}
\end{preguntabox}

\begin{preguntabox}
\section{Pregunta 2}
Las señales que transmite una antena de radio son ondas electromagnéticas. En cuanto a su tipo de vibración, estas ondas son:

\begin{enumerate}
    \item[A)] Longitudinales, porque viajan en la misma dirección en que oscilan.
    \item[B)] Longitudinales, porque no requieren medio material.
    \item[C)] Transversales, porque los campos eléctrico y magnético oscilan perpendicularmente a la dirección de propagación.
    \item[D)] Ni transversales ni longitudinales; las ondas electromagnéticas no tienen dirección de vibración.
\end{enumerate}
\end{preguntabox}

\begin{preguntabox}
\section{Pregunta 3}
El corazón de un atleta late a una frecuencia de 4 Hz durante el esfuerzo máximo. ¿Cuál es el periodo entre latidos?

\begin{enumerate}
    \item[A)] 4 s
    \item[B)] 0.4 s
    \item[C)] 0.25 s
    \item[D)] 2.5 s
\end{enumerate}
\end{preguntabox}
\newpage
\begin{preguntabox}
\section{Pregunta 4}
Un aerogenerador produce vibraciones con longitud de onda de 8 m y frecuencia de 15 Hz. ¿Con qué velocidad se propagan estas vibraciones?

\begin{enumerate}
    \item[A)] 1.875 m/s
    \item[B)] 23 m/s
    \item[C)] 120 m/s
    \item[D)] 7 m/s
\end{enumerate}
\end{preguntabox}

\begin{preguntabox}
\section{Pregunta 5}
Un tren emite un silbato que produce ondas sonoras que viajan a 400 m/s. Si el periodo de la onda es de 0.5 s, ¿cuál es su longitud de onda?

\begin{enumerate}
    \item[A)] 0.00125 m
    \item[B)] 800 m
    \item[C)] 400.5 m
    \item[D)] 200 m
\end{enumerate}
\end{preguntabox}

% -------------------------------------------------------
%  PREGUNTAS 6-10: Combinación de conceptos
% -------------------------------------------------------

\begin{preguntabox}
\section{Pregunta 6}
Un dron emite ultrasonidos que viajan a 240 m/s con una longitud de onda de 12 m. ¿Cuál es la frecuencia de la onda?

\begin{enumerate}
    \item[A)] 2 880 Hz
    \item[B)] 228 Hz
    \item[C)] 0.05 Hz
    \item[D)] 20 Hz
\end{enumerate}
\end{preguntabox}

\begin{preguntabox}
\section{Pregunta 7}
Dos olas en el océano se superponen en el mismo punto. La primera tiene una amplitud de 5 m y la segunda de 3 m. Si llegan en fase (cresta con cresta), ¿cuál es la amplitud de la ola resultante?

\begin{enumerate}
    \item[A)] 2 m
    \item[B)] 4 m
    \item[C)] 15 m
    \item[D)] 8 m
\end{enumerate}
\end{preguntabox}

\begin{preguntabox}
\section{Pregunta 8}
Los audífonos con cancelación activa de ruido generan una onda que se superpone al ruido ambiental. Si el ruido tiene amplitud de 7 unidades y la onda generada tiene amplitud de 4 unidades llegando en contrafase, ¿cuál es la amplitud resultante?

\begin{enumerate}
    \item[A)] 11 unidades
    \item[B)] 3 unidades
    \item[C)] 28 unidades
    \item[D)] 1.75 unidades
\end{enumerate}
\end{preguntabox}

\begin{preguntabox}
\section{Pregunta 9}
Una onda sonora viaja por el aire y pasa a un medio sólido donde su velocidad aumenta. ¿Cuál de las siguientes afirmaciones sobre la refracción es correcta?

\begin{enumerate}
    \item[A)] La frecuencia aumenta porque el medio es más rígido.
    \item[B)] La longitud de onda disminuye porque la velocidad aumenta.
    \item[C)] La frecuencia no cambia, pero la longitud de onda aumenta porque la velocidad aumenta.
    \item[D)] Tanto la frecuencia como la longitud de onda permanecen constantes.
\end{enumerate}
\end{preguntabox}

\begin{preguntabox}
\section{Pregunta 10}
Un médico usa ultrasonido para obtener imágenes de tejidos internos. La sonda emite ondas con $\lambda = 0.5$ mm. Una membrana celular tiene un espesor de 0.5 mm. ¿Qué fenómeno explica que el ultrasonido pueda ``rodear'' estructuras de ese tamaño?

\begin{enumerate}
    \item[A)] Reflexión, porque el ultrasonido rebota en las membranas.
    \item[B)] Refracción, porque el ultrasonido cambia de velocidad al entrar al tejido.
    \item[C)] Difracción, porque el tamaño de la estructura es comparable a $\lambda$.
    \item[D)] Interferencia, porque dos haces de ultrasonido se superponen.
\end{enumerate}
\end{preguntabox}

% -------------------------------------------------------
%  PREGUNTAS 11-15: Problemas desafiantes
% -------------------------------------------------------

\begin{preguntabox}
\section{Pregunta 11}
La señal de Wi-Fi de un router tiene una amplitud $A_0$ al salir del dispositivo. Al atravesar dos paredes, su amplitud se reduce a $A_0/2$. ¿En qué fracción queda la energía de la señal respecto a la original?

\begin{enumerate}
    \item[A)] $\dfrac{1}{2}$
    \item[B)] $\dfrac{1}{4}$
    \item[C)] $\dfrac{1}{8}$
    \item[D)] $\dfrac{1}{3}$
\end{enumerate}
\end{preguntabox}

\begin{preguntabox}
\section{Pregunta 12}
Un buque oceanográfico emite pulsos de sonar con longitud de onda de 25 m y frecuencia de 8 Hz. La señal rebota en el fondo marino y regresa 8 segundos después. ¿Cuál es la profundidad del océano en ese punto?

\begin{enumerate}
    \item[A)] 1 600 m
    \item[B)] 400 m
    \item[C)] 800 m
    \item[D)] 3 200 m
\end{enumerate}
\end{preguntabox}

\begin{preguntabox}
\section{Pregunta 13}
Una onda sísmica viaja por la corteza terrestre a 400 m/s con una frecuencia de 50 Hz. Al pasar al manto terrestre, su velocidad disminuye a 200 m/s. ¿Cuál es la longitud de onda en el manto?

\begin{enumerate}
    \item[A)] 8 m
    \item[B)] 4 m
    \item[C)] 2 m
    \item[D)] 16 m
\end{enumerate}
\end{preguntabox}
\newpage
\begin{preguntabox}
\section{Pregunta 14}
Un sensor de vibraciones en un puente registra 18 ciclos completos en 3 segundos. Si la velocidad de propagación de la vibración en la estructura es de 360 m/s, ¿cuál es la longitud de onda?

\begin{enumerate}
    \item[A)] 6 m
    \item[B)] 120 m
    \item[C)] 20 m
    \item[D)] 60 m
\end{enumerate}
\end{preguntabox}

\begin{preguntabox}
\section{Pregunta 15}
En un sistema de radar meteorológico, las microondas viajan a 680 m/s en cierto medio y su eco regresa 0.5 s después. Un ingeniero modifica el equipo para que emita la misma onda pero con el doble de longitud de onda. ¿A qué nueva frecuencia debe operar el radar?

\begin{enumerate}
    \item[A)] $2f$
    \item[B)] $4f$
    \item[C)] $\dfrac{f}{4}$
    \item[D)] $\dfrac{f}{2}$
\end{enumerate}
\end{preguntabox}

% -------------------------------------------------------
%  RESPUESTAS Y RETROALIMENTACIÓN
% -------------------------------------------------------

\newpage
\onecolumn
\section*{RESPUESTAS Y RETROALIMENTACIÓN}

\begin{respuestabox}
\textbf{Pregunta 1}\\
\textcolor{verdecorrecto}{\textbf{Respuesta correcta: B) Longitudinales, porque son ondas de presión.}}

\textbf{Justificación:}\\
Las ondas de sonar son ondas sonoras: perturbaciones de presión que se propagan por compresión y rarefacción del agua en la misma dirección de propagación. (Lección: Sección 1 — Clasificación por la dirección de vibración)

\textbf{Respuestas incorrectas:}
\begin{itemize}
    \item \textbf{A):} El medio (líquido o sólido) no determina el tipo de onda; lo determina la dirección de la perturbación respecto a la propagación.
    \item \textbf{C):} La dirección de movimiento del submarino no tiene relación con el tipo de onda emitida.
    \item \textbf{D):} Las ondas electromagnéticas no requieren medio, pero el sonar sí lo requiere; no puede propagarse en el vacío.
\end{itemize}
\end{respuestabox}

\begin{respuestabox}
\textbf{Pregunta 2}\\
\textcolor{verdecorrecto}{\textbf{Respuesta correcta: C) Transversales, porque los campos oscilan perpendicularmente a la propagación.}}

\textbf{Justificación:}\\
En una onda electromagnética, los campos eléctrico y magnético oscilan en planos perpendiculares a la dirección de propagación, lo que las define como transversales. (Lección: Sección 1 — Onda transversal)

\textbf{Respuestas incorrectas:}
\begin{itemize}
    \item \textbf{A) y B):} Confunden la característica de no requerir medio con la dirección de vibración; son aspectos distintos.
    \item \textbf{D):} Todas las ondas tienen una dirección de vibración definida; en las EM es perpendicular a la propagación.
\end{itemize}
\end{respuestabox}

\begin{respuestabox}
\textbf{Pregunta 3}\\
\textcolor{verdecorrecto}{\textbf{Respuesta correcta: C) 0.25 s}}

\textbf{Fórmula utilizada:} $T = \dfrac{1}{f}$ (Lección: Sección 1 — Periodo)

\textbf{Sustitución:}
$$T = \frac{1}{4 \text{ Hz}} = 0.25 \text{ s}$$

\textbf{Respuestas incorrectas:}
\begin{itemize}
    \item \textbf{A) 4 s:} Copia el valor de $f$ directamente como si fuera $T$.
    \item \textbf{B) 0.4 s:} Calcula $\frac{1}{2.5}$ en lugar de $\frac{1}{4}$; error aritmético.
    \item \textbf{D) 2.5 s:} Calcula $\frac{10}{4}$; introduce un factor de 10 sin justificación.
\end{itemize}
\end{respuestabox}
\newpage
\begin{respuestabox}
\textbf{Pregunta 4}\\
\textcolor{verdecorrecto}{\textbf{Respuesta correcta: C) 120 m/s}}

\textbf{Fórmula utilizada:} $v = \lambda \cdot f$ (Lección: Sección 1 — Velocidad de propagación)

\textbf{Sustitución:}
$$v = 8 \text{ m} \times 15 \text{ Hz} = 120 \text{ m/s}$$

\textbf{Respuestas incorrectas:}
\begin{itemize}
    \item \textbf{A) 1.875 m/s:} Divide $\lambda \div f = 8 \div 15$; invierte la operación correcta $\lambda \div f$ para obtener $T$, pero lo confunde con $v$.
    \item \textbf{B) 23 m/s:} Suma $\lambda + f = 8 + 15$; confunde la operación.
    \item \textbf{D) 7 m/s:} Resta $15 - 8$; confunde la operación.
\end{itemize}
\end{respuestabox}

\begin{respuestabox}
\textbf{Pregunta 5}\\
\textcolor{verdecorrecto}{\textbf{Respuesta correcta: D) 200 m}}

\textbf{Fórmula utilizada:} $\lambda = v \cdot T$ (Lección: Sección 1 — Velocidad de propagación)

\textbf{Sustitución:}
$$\lambda = 400 \text{ m/s} \times 0.5 \text{ s} = 200 \text{ m}$$

\textbf{Respuestas incorrectas:}
\begin{itemize}
    \item \textbf{A) 0.00125 m:} Calcula $T \div v = 0.5 \div 400$; invierte el cociente.
    \item \textbf{B) 800 m:} Calcula $v \div T = 400 \div 0.5 = 800$; confunde $\lambda = v \cdot T$ con $\lambda = v \div T$.
    \item \textbf{C) 400.5 m:} Suma $v + T = 400 + 0.5$; confunde la operación.
\end{itemize}
\end{respuestabox}

\begin{respuestabox}
\textbf{Pregunta 6}\\
\textcolor{verdecorrecto}{\textbf{Respuesta correcta: D) 20 Hz}}

\textbf{Fórmula utilizada:} $f = \dfrac{v}{\lambda}$ (Lección: Sección 1 — Velocidad de propagación)

\textbf{Sustitución:}
$$f = \frac{240 \text{ m/s}}{12 \text{ m}} = 20 \text{ Hz}$$

\textbf{Respuestas incorrectas:}
\begin{itemize}
    \item \textbf{A) 2 880 Hz:} Multiplica $v \times \lambda = 240 \times 12$; invierte la operación.
    \item \textbf{B) 228 Hz:} Resta $240 - 12$; confunde la operación.
    \item \textbf{C) 0.05 Hz:} Calcula $\lambda \div v = 12 \div 240$; invierte el cociente.
\end{itemize}
\end{respuestabox}

\begin{respuestabox}
\textbf{Pregunta 7}\\
\textcolor{verdecorrecto}{\textbf{Respuesta correcta: D) 8 m}}

\textbf{Fórmula utilizada:} $A_{\text{total}} = A_1 + A_2$ (Lección: Sección 2 — Interferencia constructiva)

\textbf{Sustitución:}
$$A_{\text{total}} = 5 \text{ m} + 3 \text{ m} = 8 \text{ m}$$

\textbf{Respuestas incorrectas:}
\begin{itemize}
    \item \textbf{A) 2 m:} Aplica $|A_1 - A_2|$; usa la fórmula de interferencia \textit{destructiva} en lugar de constructiva.
    \item \textbf{B) 4 m:} Promedia las amplitudes $(5+3)/2 = 4$; confunde promedio con superposición.
    \item \textbf{C) 15 m:} Multiplica $A_1 \times A_2 = 5 \times 3$; confunde la operación.
\end{itemize}
\end{respuestabox}

\begin{respuestabox}
\textbf{Pregunta 8}\\
\textcolor{verdecorrecto}{\textbf{Respuesta correcta: B) 3 unidades}}

\textbf{Fórmula utilizada:} $A_{\text{total}} = |A_1 - A_2|$ (Lección: Sección 2 — Interferencia destructiva)

\textbf{Sustitución:}
$$A_{\text{total}} = |7 - 4| = 3 \text{ unidades}$$

\textbf{Respuestas incorrectas:}
\begin{itemize}
    \item \textbf{A) 11 unidades:} Suma las amplitudes $7 + 4$; aplica interferencia constructiva en lugar de destructiva.
    \item \textbf{C) 28 unidades:} Multiplica $7 \times 4$; confunde la operación.
    \item \textbf{D) 1.75 unidades:} Calcula $7 \div 4$; confunde la operación.
\end{itemize}
\end{respuestabox}

\begin{respuestabox}
\textbf{Pregunta 9}\\
\textcolor{verdecorrecto}{\textbf{Respuesta correcta: C) La frecuencia no cambia, pero la longitud de onda aumenta.}}

\textbf{Justificación:}\\
Al refractarse, la frecuencia siempre se conserva (es una propiedad de la fuente, no del medio). Como $v = \lambda \cdot f$ y $v$ aumenta con $f$ constante, entonces $\lambda$ debe aumentar proporcionalmente. (Lección: Sección 2 — Refracción)

\textbf{Respuestas incorrectas:}
\begin{itemize}
    \item \textbf{A):} La frecuencia nunca cambia al cambiar de medio; es la longitud de onda la que se adapta.
    \item \textbf{B):} Invierte la relación correcta: si $v$ aumenta y $f$ es constante, $\lambda$ \textit{aumenta}, no disminuye.
    \item \textbf{D):} La longitud de onda sí cambia; solo la frecuencia permanece constante.
\end{itemize}
\end{respuestabox}

\begin{respuestabox}
\textbf{Pregunta 10}\\
\textcolor{verdecorrecto}{\textbf{Respuesta correcta: C) Difracción, porque el tamaño de la estructura es comparable a $\lambda$.}}

\textbf{Justificación:}\\
La difracción es significativa cuando $d \approx \lambda$. Aquí, el espesor de la membrana ($d = 0.5$ mm) es exactamente igual a la longitud de onda del ultrasonido ($\lambda = 0.5$ mm), lo que permite que la onda se doble alrededor de la estructura. (Lección: Sección 2 — Difracción)

\textbf{Respuestas incorrectas:}
\begin{itemize}
    \item \textbf{A):} La reflexión ocurre en interfaces, no explica el rodeo de estructuras pequeñas.
    \item \textbf{B):} La refracción explica el cambio de dirección al cambiar de medio, no el rodeo de obstáculos.
    \item \textbf{D):} La interferencia requiere la superposición de dos o más ondas; no explica el comportamiento alrededor de una estructura.
\end{itemize}
\end{respuestabox}

\begin{respuestabox}
\textbf{Pregunta 11}\\
\textcolor{verdecorrecto}{\textbf{Respuesta correcta: B) $\dfrac{1}{4}$}}

\textbf{Fórmula utilizada:} $E \propto A^2$ (Lección: Sección 2 — Amortiguación)

Si $A' = \dfrac{A_0}{2}$:
$$E' \propto \left(\frac{A_0}{2}\right)^2 = \frac{A_0^2}{4} \quad \Rightarrow \quad E' = \frac{E_0}{4}$$

\textbf{Respuestas incorrectas:}
\begin{itemize}
    \item \textbf{A) $\frac{1}{2}$:} Error de pensar que la energía es proporcional a $A$ (lineal) en lugar de $A^2$ (cuadrática).
    \item \textbf{C) $\frac{1}{8}$:} Aplica $A^3$ en lugar de $A^2$; error en el exponente.
    \item \textbf{D) $\frac{1}{3}$:} No corresponde a ninguna relación física correcta entre amplitud y energía.
\end{itemize}
\end{respuestabox}

\begin{respuestabox}
\textbf{Pregunta 12}\\
\textcolor{verdecorrecto}{\textbf{Respuesta correcta: C) 800 m}}

\textbf{Paso 1 — Calcular velocidad:} $v = \lambda \cdot f$ (Lección: Sección 1)
$$v = 25 \text{ m} \times 8 \text{ Hz} = 200 \text{ m/s}$$

\textbf{Paso 2 — Calcular profundidad} (señal recorre ida y vuelta en 8 s):
$$d = v \cdot \frac{t}{2} = 200 \text{ m/s} \times \frac{8 \text{ s}}{2} = 200 \times 4 = 800 \text{ m}$$

\textbf{Respuestas incorrectas:}
\begin{itemize}
    \item \textbf{A) 1 600 m:} Olvida dividir el tiempo entre 2; calcula $200 \times 8$ sin considerar el viaje de regreso.
    \item \textbf{B) 400 m:} Usa solo $\lambda$ como velocidad: $25 \times 4 \times 4$; ignora la fórmula $v = \lambda f$.
    \item \textbf{D) 3 200 m:} Multiplica $v \times t$ sin dividir y además duplica: $200 \times 8 \times 2$.
\end{itemize}
\end{respuestabox}

\begin{respuestabox}
\textbf{Pregunta 13}\\
\textcolor{verdecorrecto}{\textbf{Respuesta correcta: B) 4 m}}

\textbf{Justificación:}\\
Al refractarse, la frecuencia no cambia ($f = 50$ Hz). En el manto: (Lección: Sección 2 — Refracción)

\textbf{Sustitución:}
$$\lambda_2 = \frac{v_2}{f} = \frac{200 \text{ m/s}}{50 \text{ Hz}} = 4 \text{ m}$$

\textbf{Respuestas incorrectas:}
\begin{itemize}
    \item \textbf{A) 8 m:} Usa la velocidad de la corteza ($\lambda_1 = 400/50 = 8$ m); no aplica el cambio de velocidad al cruzar al manto.
    \item \textbf{C) 2 m:} Calcula $v_2 \div v_1 = 200 \div 400 = 0.5$, luego $\lambda_1 \times 0.5 / f$; mezcla los pasos incorrectamente.
    \item \textbf{D) 16 m:} Suma $v_1 + v_2 = 600$, luego divide entre $f$; confunde la operación.
\end{itemize}
\end{respuestabox}

\begin{respuestabox}
\textbf{Pregunta 14}\\
\textcolor{verdecorrecto}{\textbf{Respuesta correcta: D) 60 m}}

\textbf{Paso 1 — Calcular frecuencia:} (Lección: Sección 1)
$$f = \frac{18 \text{ ciclos}}{3 \text{ s}} = 6 \text{ Hz}$$

\textbf{Paso 2 — Calcular longitud de onda:}
$$\lambda = \frac{v}{f} = \frac{360 \text{ m/s}}{6 \text{ Hz}} = 60 \text{ m}$$

\textbf{Respuestas incorrectas:}
\begin{itemize}
    \item \textbf{A) 6 m:} Solo calcula $f = 6$ Hz y lo reporta como $\lambda$; omite el segundo paso.
    \item \textbf{B) 120 m:} Calcula $v \times T = 360 \times (3/18) \times 2$; introduce un factor extra.
    \item \textbf{C) 20 m:} Calcula $v \div$ (número de ciclos) $= 360 \div 18 = 20$; omite convertir ciclos a frecuencia usando el tiempo.
\end{itemize}
\end{respuestabox}
\newpage
\begin{respuestabox}
\textbf{Pregunta 15}\\
\textcolor{verdecorrecto}{\textbf{Respuesta correcta: D) $\dfrac{f}{2}$}}

\textbf{Justificación:}\\
La velocidad de propagación en el medio no cambia por modificar el equipo emisor. De $v = \lambda \cdot f$ con $v$ constante: (Lección: Sección 1 — Nota sobre relación $\lambda$-$f$)

Si $\lambda' = 2\lambda$:
$$f' = \frac{v}{\lambda'} = \frac{v}{2\lambda} = \frac{f}{2}$$

Al duplicar la longitud de onda, la frecuencia se reduce a la mitad. (El dato de la distancia al eco y la velocidad 680 m/s son contexto; no afectan la relación $\lambda$-$f$.)

\textbf{Respuestas incorrectas:}
\begin{itemize}
    \item \textbf{A) $2f$:} Confunde proporcionalidad directa con inversa; si $\lambda$ sube, $f$ baja, no sube.
    \item \textbf{B) $4f$:} Aplica una relación cuadrática inexistente entre $\lambda$ y $f$.
    \item \textbf{C) $\frac{f}{4}$:} Aplica la relación de energía ($E \propto A^2$) a la frecuencia; confunde conceptos.
\end{itemize}
\end{respuestabox}